\section{Discussion}\label{discussion}

\subsection{Compliance-by-construction is reachable on consumer hardware}\label{compliance-by-construction-is-reachable-on-consumer-hardware}

The pipeline described here demonstrates that PhysioNet-compliant
fine-tuning of a research-grade biomedical classifier is reachable
on a single consumer Apple Silicon laptop without cloud GPU
access, institutional compute budget, or BAA negotiation. The
\textasciitilde15-hour wall-clock per training run is slow relative to A100
cloud GPUs (\textasciitilde3--5× slower), but the cost is one-time per
checkpoint and the data-residency story is dramatically simpler.
For credentialed researchers who do not have institutional cloud
infrastructure pre-approved for credentialed clinical data, the
local-MPS path is the path of least friction.

The same code path supports CUDA backends; researchers with
appropriate cloud-compute authorization can swap
\texttt{torch.device("mps")} for \texttt{torch.device("cuda")} and gain the
wall-clock speedup. The compliance posture in that case becomes a
function of the cloud environment's BAA and data-residency
review, not the training code itself.

\subsection{In-distribution vs.~out-of-distribution gap}\label{in-distribution-vs.-out-of-distribution-gap}

The 99\%-validation / 40--50\%-Synthea gap is the most
deployment-relevant single finding from this work. It reflects a
common pattern in biomedical NLP that in-distribution validation
numbers can dramatically overstate real-world performance when
the deployment distribution differs in surface form or cohort
composition from the training distribution. Reporting \emph{only}
in-distribution validation numbers is therefore misleading; we
report both, and we argue that any clinical deployment claim made
on the basis of in-distribution validation accuracy alone should
be regarded with skepticism.

Synthea is not the only out-of-distribution test that could be
run. Other reasonable choices include: cross-institution MIMIC
holdouts where available, pediatric or non-English clinical
corpora, or institution-specific notes from the deploying
hospital. Any such test would surface a gap of similar character;
our headline number is the one that the companion paper
\autocite{FungPhantomCodes2026} uses to anchor its non-LLM baseline.

\subsection{Limitations}\label{limitations}

\begin{itemize}
\tightlist
\item
  \textbf{Top-50 vocabulary cap by construction.} The classifier
  cannot predict any code outside the 50 most-frequent training
  codes. This is a deliberate trade-off (signal density over
  coverage), but it bounds applicability to deployments where
  long-tail codes are not central. A v2 variant with a larger
  head or with hierarchical decomposition is straightforward to
  build and is on the roadmap.
\item
  \textbf{Single random seed for v1.} Reported numbers come from a
  single seed. A proper sensitivity analysis (5+ seeds with
  confidence intervals) is on the v2 roadmap; for v1 the
  single-seed result is reported alongside this explicit
  acknowledgment.
\item
  \textbf{D1/D2 sequence-length truncation.} Setting
  \texttt{max\_seq\_length=128} for hardware feasibility truncates the
  longest D1\_full and D2\_no\_code FHIR-JSON inputs. The
  D3\_text\_only and D4\_abbreviated headline inputs (10--50
  tokens) are unaffected; the truncation matters most for the
  D1/D2 modes that include serialized JSON payloads.
\item
  \textbf{MIMIC-IV training cohort scope.} ACCESS-scope conditions
  are weighted toward cardiometabolic disease; performance on
  conditions outside that scope is not directly addressable from
  this trained model.
\item
  \textbf{No comparison against larger biomedical encoders.} BiomedLM,
  GatorTron, BioMistral, and Med-PaLM (via parameter-efficient
  fine-tuning) are plausible competitors; we deferred those
  comparisons to a v2 ablation arm rather than crowd v1.
\end{itemize}

\subsection{Future directions}\label{future-directions}

\begin{itemize}
\tightlist
\item
  \textbf{Parameter-efficient fine-tuning} of larger biomedical encoders
  via LoRA \autocite{Hu2022} or QLoRA \autocite{Dettmers2023} would close the gap
  between this 110M-parameter classifier and larger biomedical
  models without exceeding our hardware constraints.
\item
  \textbf{Vocabulary expansion to LOINC and RxNorm} for laboratory
  observations and medications, using the same architecture and
  training recipe.
\item
  \textbf{Multi-seed sensitivity analysis} with confidence intervals
  reported alongside point estimates.
\item
  \textbf{Cross-institution out-of-distribution testing} where
  credentialed cohorts from additional institutions are
  available.
\end{itemize}
